\section{Discussion and Merge Boundary}

\subsection{What the current result establishes}

The retained route demonstrates that low-precision backward should be treated
as a represented operator, not as BF16 backward followed by casts.  The fixed
powers in Equations~\ref{eq:fp8-operands}--\ref{eq:represented-dk} remove
general rescaling from the hot path; P and dS are quantized at production;
and adaptive Q/K metadata is consumed by tensor scale pages and gradient
epilogues without adding an MMA.

The learned-projection result adds the systems half of that argument.  NVFP4
tensor work is useful only if its operand preparation and output layouts are
part of surrounding producers.  The unified epilogue emits one Q/K payload
for both directions, transposed MXFP4 V, and optional BF16 Q/K/V.  It thereby
avoids multiple BF16 rereads, standalone transposes, and fake-quantization
launches, producing a much larger systems gain than another local backward
barrier adjustment.

The final pure-FP4 specialization applies that rule more aggressively.  It
quantizes each Q/K projection fragment once at the fixed pure scale and fans
the same codes into all six consumer layouts.  This removes 69.6--97.2
$\mu$s from two projection-plus-backward chains.  The result also separates
the remaining limits cleanly: pure operand generation is no longer the main
setup penalty, yet the mixed-dP hybrid remains faster because block maxima,
dual dS scale publication, and score fanout are exposed inside backward.
Changing those internal pages from MXFP4 K32/E8M0 to NVFP4 K16/E4M3 does not
recover the gap.  The more accurate NVFP4-dS variant is slower than BF16, and
the combined NVFP4 diagnostic is non-finite.

The lifetime-safe adaptive mixed-dP result is a narrower positive outcome.
Once the dP tail word and adaptive score page have separate lifetimes, the
route is stable at the same zero-spill register limit.  It wins at H24 but not
H8 and remains in the approximately 0.983 dQ/dK accuracy class.  This confirms
that adaptive scaling and mixed operand precision are independent choices;
fixing the former does not erase error introduced by the latter.

The upstream dO result supports the same rule from the other side of the
graph.  Row and column MXFP4 representations and the softmax delta should be
published by one producer, not reconstructed through three kernels immediately
before backward.  This halves that setup boundary without changing any
backward arithmetic.

The composed MFU result tests the rule at a model-facing boundary.  Keeping
the FP4 FA4 forward and FP4+FP8 backward together raises useful algorithmic
MFU on all five measured shapes.  The gain is modest at H16 and largest at
wide H64 projection widths, where projection-native NVFP4 work amortizes
packing and publication.  The faster activation-only stacked handoff is not
silently used for full training: the latter retains BF16 QKV publication so
the QKV weight-gradient GEMM remains present.

\subsection{Why backward does not reproduce the forward gain directly}

The companion forward report uses two matrix products and two major TMEM
accumulators, then attacks the serial score-to-P publication path.  Backward
has five products, persistent dK/dV state, two dS orientations, and a
cross-owner dQ reduction.  Its critical path is therefore not just
exponential evaluation.  Approximating $2^x$ more aggressively cannot remove
the score TMEM fanout, dP centering, dS publication, or dQ reduction.

The current backward also chooses E4M3 P for its accuracy-oriented route,
whereas the forward report's retained full-FP4 policy uses MXFP4 P/V and a
direct E2M1 code map.  Those are different represented operators.  A local
represented-P replay saves 0.76\%, but materializing forward P for backward
would write and reread 272 MiB at S8192/H8.  The profitable analogue of
forward reuse must therefore remain on chip or change the consumer topology;
an HBM cache is not a useful bridge.

\subsection{What the six integration tests rule out}

The tests narrow the next breakthrough substantially.  It is not a second Q/K
format, a standalone V quantizer, three separate dO setup kernels, a global P
cache, or another barrier.  Per-head polling inside the projection K loop is
enough for a small H8 win, and double-buffering the projection's private TMEM
removes its MMA--epilogue bubble.  It is not enough at wider heads.  The
remaining high-value boundary is the global dQ reduction: projection backward
must consume the reduction tile before, or instead of, BF16 global
materialization.  Observing the tile immediately after publication still
pays the global rendezvous and leaves a bounded-consumer tail.

The compact NVFP4 final-tile handoff sharpens this boundary.  Its isolated
projection, including delayed-scale conversion, beats BF16 at both measured
widths and reaches approximately 0.991 cosine.  The composed chain wins 3.7\%
at H24 but loses 1.9\% at H8 because it still waits for and rereads the private
BF16 reduction.  Format choice can accelerate the consumer; it cannot by
itself remove the producer topology that defines the ceiling.

The implemented two-lane hierarchy further rules out a weaker interpretation
of that target.  Projection can consume two BF16 owner sums directly and avoid
ever constructing completed dQ, but it is slower because BF16 partial traffic
is unchanged and the consumer must duplicate the weight MMA or fold the lanes.
The useful breakthrough must therefore combine projection integration with
pre-publication owner grouping or a compact FP8/FP4 partial representation.
Moving the final addition alone does not approach the no-publication timing
diagnostic.

The stacked-NVFP4 version reaches the same conclusion through a different
consumer.  Folding two dQ lanes, inverse RoPE, and local quantization in one
K128 producer is accurate, but its readiness signals perturb the attention
schedule enough to lose.  Removing those signals makes the one-lane stacked
layout a small activation-only winner, but also restores completed BF16 dQ.
It is a useful publication optimization, not the sought reduction-topology
breakthrough.

Storage reclamation also has a hard hardware boundary.  Projection scratch
can be unioned because Q/K code staging and V BF16-pair staging have disjoint
lifetimes.  The backward mixed-dP row cannot simply shrink from 128 to 96
bytes: the grouped SM100 depth transaction does not retire with that geometry.
The final 32 bytes are logically padding but are not currently reclaimable by
a legal drop-in descriptor.

\subsection{Accuracy and training limits}

The reported cosine values cover synthetic calibrated, unit-scale, and
sparse-outlier regimes.  The new 500-step teacher/student experiment adds a
block-level optimization trajectory with refreshed weights and complete
projection weight gradients.  Across three seeds, low precision retains
89.7--99.6\% of BF16's absolute validation-loss improvement.  This does not
establish language-model convergence, final perplexity, optimizer-state
stability over many layers, or sensitivity to realistic activation tails.
The learned NVFP4 projection itself has Q/K/V cosine near 0.991, while
attention backward has a separate approximately 0.991 dQ/dK cosine.  Native
forward output cosine is 0.9907--0.9913 on the two unified-projection checks.
The compact NVFP4 dQ projection contribution has approximately 0.991 cosine
against the BF16 projection; FP8 is approximately 0.9993 and MXFP4
approximately 0.9832--0.9834 on the same boundary.
Their composed error is not the product of those scalars.  The trajectory
experiment shows that reducible error falls, but the FP4 forward contributes
a native teacher floor of 0.093--0.098 relative MSE and the final shared-BF16
evaluation remains 25.5\% above the BF16 student's loss on average.  The raw
loss gap must remain visible beside the floor-adjusted convergence rate.

Robust per-head clipping deliberately sacrifices extreme values to protect
the bulk distribution.  Per-token or blockwise Q/K scales could retain more
outliers, but scales varying along the $dQ/dK$ reduction dimension require
additional instruction scale pages and TMEM movement.  Quantization-aware
training is the natural larger validation framework
\cite{zhang2026attnqat}.

\subsection{Hardware and shape scope}

The active specialization is causal, batch one, sequence divisible by 128,
$D_{QK}=192$, and $D_V=128$ on data-center Blackwell.  It supports the
measured S4096, S8192, and S16384 families and selected head counts.  The
projection kernel requires matrix dimensions aligned to its NVFP4 tiling.
GQA, variable-length batches, dropout, arbitrary head dimensions, and a
portable non-TMEM implementation are outside this report.

The 512-column alias schedule, \tcgen{} scale pages, and two-CTA matrix shapes
are SM100-specific \cite{nvidia2026ptx}.  A device with another score bank,
different scaled-MMA K granularity, or cheaper cross-cluster reduction could
choose a different ownership plan.

\subsection{What belongs to prior structure}

FlashAttention supplies the recomputation and IO-aware backward formulation
\cite{dao2022flashattention}.  FA4 supplies the hardware-asymmetric,
warp-specialized scheduling framework \cite{zadouri2026flashattention4}.
ThunderKittens supplies the optimized V382 control and the tile abstractions
\cite{thunderkittens}.  HAO's forward FP4 work is a structural reference for
block-scaled tensor operations, TMEM lifetime reuse, and producer-side
quantization \cite{hao2026fp4flashattention}.  This report's contribution is
the backward precision factorization, adaptive Q/K contract, retained TMEM
handoffs, BF16 dQ endpoint, unified learned-QKV publication, fused upstream dO
contract, and measured projection-gradient boundary built on those
foundations.

\subsection{Should this merge with the forward report?}

Not yet as a single method section.  The reports should remain separate while
the following are unresolved:

\begin{itemize}
  \item the forward and backward P representations differ, and the measured
  HBM P cache is not viable;
  \item projection-scale production is still an explicit integration
  contract;
  \item composed forward-plus-backward accuracy and a synthetic block-level
  trajectory are measured, but stock-model loss and perplexity are not;
  \item the pure-FP4 backward route is validated but experimental, while the
  retained backward is the more accurate FP4+FP8 route;
  \item the fused dO producer and all-MX consumer compose successfully, but
  that endpoint is not the selected performance/accuracy frontier.
\end{itemize}

A later umbrella manuscript should merge the shared introduction, Blackwell
format background, notation, and end-to-end evaluation.  It should keep
separate forward and backward architecture/method sections, then add one
systems section for projection epilogues and one composed training-results
section.  This avoids duplicating background without implying that the two
kernels currently implement the same probability operator.
